\section{Adaptive quartet conditioning in polynomial space}
\label{sec:quartet}

The incoming adaptive-contraction packet~\cite{incomingadaptive} already uses balanced binary
separator conditioning. The new step in this section is to choose the
binary pairing of four fixed blocks from the actual exact-weight list
sizes. This choice is repeated inside the blocks. It changes the
worst-case coefficient of the paid child-call count; it is not merely a
retuning of the taper schedule.

We use the common normalized contraction interface and its imported
asymmetric packing, slope-four Cauchy separation, precision, and child
construction results. For one exact pair-budget array, let
$\mathcal I_{xy}$ be the list on directed group pair $(x,y)$, discard an
array with an empty list, and put
\[
 h_{xy}=\log_2|\mathcal I_{xy}|,\qquad
 \mathcal K(M,u)=\kappa\frac{M}{\sqrt u},\qquad
 \kappa=\frac{\pi\sqrt2}{\ln2}.
\]
Lemma~\ref{lem:norm} implies, for every fixed nonnegative
coefficient vector $z$ on the directed group pairs,
\begin{equation}\label{eq:quartet-entropy}
 \sum_{x\ne y}z_{xy}h_{xy}
 \le\left(\frac{\|z\|_2}{r}+o(1)\right)\mathcal K(M,u).
\end{equation}
All coefficient vectors used below belong to a finite family independent
of $s$, since every alphabet size is fixed. Thus the bound is uniform
even though our algorithm selects a vector indirectly from the list sizes.
This avoids an unjustified uniformity inference for an arbitrary
continuously varying $z$.

\subsection{The four-block identity and its norm bound}

Partition a set of groups into four fixed nonempty blocks
$B_1,B_2,B_3,B_4$. Let
\[
 H_{ab}=\sum_{x\in B_a,\,y\in B_b}(h_{xy}+h_{yx})
 \quad (a<b)
\]
be the logarithm of the product of list sizes crossing a pair of blocks.
There are three perfect matchings of the four blocks. They partition the
six unordered block-pair classes into three complementary pairs. In the
$j$th such pair write its two weights as $a_j\ge b_j$.

For a pairing $(B_a\cup B_b)\mid(B_c\cup B_d)$, condition first on
all directed indices between the two sides, then condition inside each
side between its two constituent blocks. The largest external index
weight on a path ending at one of the four block oracles is the total
cross-block weight minus $\min(H_{ab},H_{cd})$. Consequently the best
pairing has exact external path bound
\begin{equation}\label{eq:quartet-F}
 F_4(H)=\sum_{j=1}^3(a_j+b_j)-\max_j b_j.
\end{equation}
Its choice is computable without real logarithms: for each matching take
the smaller of its two integer products of list sizes, and maximize that
integer. There are just three candidates.

\begin{lemma}[Adaptive quartet inequality]\label{lem:quartet}
For every nonnegative six-edge weighting $H$,
\[
 F_4(H)\le\sum_j a_j+\frac23\sum_j b_j
 \le\sqrt{\frac{13}{3}}\left(\sum_{a<b}H_{ab}^2\right)^{1/2}.
\]
Both inequalities can hold with equality. For four singleton groups,
there is a directed-edge coefficient vector of squared norm $26/3$
whose scalar product with $h$ bounds the maximum path weight of the
selected separator tree.
\end{lemma}
\begin{proof}
Since $\max_j b_j\ge\frac13\sum_j b_j$, the first inequality follows
from~\eqref{eq:quartet-F}. Cauchy--Schwarz gives the second because
$3\cdot1^2+3\cdot(2/3)^2=13/3$. Equality occurs when the three $a_j$
are all $3t$ and the three $b_j$ are all $2t$, for any $t\ge0$.
Each undirected singleton-pair weight is the sum of two directed weights;
the coefficient pattern therefore has squared norm $2(13/3)=26/3$.
\end{proof}

The lemma is sharp for this local family of three balanced binary trees.
It does not claim that every possible contraction algorithm has this
worst-case coefficient.

For unequal block sizes $k_a=|B_a|$, a weighted version is useful. Choose
any fixed probabilities $p_1,p_2,p_3\ge0$ with sum one. Equation
\eqref{eq:quartet-F} implies
\[
 F_4(H)\le\sum_j a_j+\sum_j(1-p_j)b_j.
\]
The complementary block-pair classes in matching $j$ contain respectively
$c_j$ and $d_j$ directed pairs, each of the form $2k_ak_b$. Regardless of
which class has the larger actual weight, the squared coefficient norm
is at most
\begin{equation}\label{eq:quartet-unequal}
 R(k,p)=\sum_{j=1}^3\left[
 \max(c_j,d_j)+\min(c_j,d_j)(1-p_j)^2\right].
\end{equation}
The probabilities are only an analysis device; the algorithm still
chooses the matching by its exact score in~\eqref{eq:quartet-F}.

\subsection{Recursive quartets and all child calls}

Apply the same rule inside each of the four blocks. Blocks of size two
use their unique binary split; singletons are leaves. Let $g_r$ bound
the squared norm of one coefficient vector that upper-bounds the largest
root-to-leaf index weight of the resulting tree on $r$ groups. If the
four block sizes are $k_1,\ldots,k_4$, then
\begin{equation}\label{eq:quartet-g-recurrence}
 g_r\le R(k,p)+\max_a g_{k_a},\qquad g_1=0,\quad g_2=2.
\end{equation}
Indeed, the external maximum is bounded by the quartet vector. The
largest internal path cost occurs in one block; choose its inductively
constructed vector. Internal and external supports are disjoint, so
squared norms add. The total cost is bounded by the sum of the two
maxima, even if they occur at different leaves. This possible
looseness is harmless for an upper bound.

For four equal blocks of size $k$, use $p_j=1/3$. Then
\[
 g_{4k}\le\frac{26}{3}k^2+g_k.
\]
For $r=4^j$ this yields
\begin{equation}\label{eq:quartet-powers}
 g_r=\frac{26}{45}(r^2-1),
\end{equation}
where equality denotes our certificate, not an optimality assertion.
The fixed balanced binary-tree certificate is $2(r^2-1)/3$; its ratio
to~\eqref{eq:quartet-powers} is exactly $15/13$. Thus the new square-root
coefficient is smaller by $\sqrt{13/15}$ for these alphabets.

The schedule below needs only the following exact ledger:
\begin{center}
\begin{tabular}{c|c|c|c|c}
$r$ & four block sizes & $(p_1,p_2,p_3)$ & $g_r$ & fixed binary $v(r)$\\\hline
4 & $(1,1,1,1)$ & $(1/3,1/3,1/3)$ & $26/3$ & $10$\\
6 & $(1,1,2,2)$ & $(0,1/2,1/2)$ & $22$ & $24$\\
8 & $(2,2,2,2)$ & $(1/3,1/3,1/3)$ & $110/3$ & $42$\\
16 & $(4,4,4,4)$ & $(1/3,1/3,1/3)$ & $442/3$ & $170$
\end{tabular}
\end{center}
For $r=6$, the complementary directed cardinalities are $(2,8)$,
$(4,4)$, and $(4,4)$. Their contribution in
\eqref{eq:quartet-unequal} is $20$, and the largest internal block
contributes $g_2=2$.

At an internal node of a separator tree, enumerate the directed indices
between its two children, recursively evaluate each conditional sum,
multiply the two answers, and accumulate. At a singleton, evaluate the
group factor by its genuine recursive tensor contraction. Finite
distributivity proves the value is exactly the original mode sum over
any commutative coefficient ring. If $E_v$ is the union of cuts on the
path to group $v$, its factor is called precisely
$\prod_{e\in E_v}|\mathcal I_e|$ times. Hence every child call is paid,
and~\eqref{eq:quartet-entropy} gives the local recurrence
\begin{equation}\label{eq:quartet-local}
 T_d\le\operatorname{poly}(s)
 2^{(\sqrt{g_{r_d}}/r_d+o(1))\mathcal K(M_d,u_d)}T_{d+1}.
\end{equation}
All loops, ranks, exact counts, current assignments, and ring accumulators
have polynomial descriptions. Depth-first traversal uses polynomial
space. There are no cached exponential tables and no free group-factor
oracle in this bound.

\subsection{A strict rational schedule}

Let $\lambda=\lceil\log_2s\rceil$. Use
\begin{equation}\label{eq:quartet-schedule}
\begin{array}{c|ccccc}
 d&0&1&2&3&4+k\ (k\ge0)\\\hline
 r_d&16&16&8&6&4\\
 \eta_d&1249/100000&1249/100000&26/3125&107/20000&(27/6400)(5/8)^k.
\end{array}
\end{equation}
Put $u_d=\lfloor\eta_d\lambda\rfloor$, use the additive buffer
$b=\lceil6\log_2\lambda\rceil$, set
$a_d=\sum_{j\le d}(u_j+b)$ with $a_{-1}=0$, and stop the network recursion at
$M\le H=\lceil s/\lambda\rceil$. The cumulative leading ranges are
\[
 A_0=\frac{1249}{100000},\quad A_1=\frac{1249}{50000},\quad
 A_2=\frac{333}{10000},\quad A_3=\frac{773}{20000},\quad
 A_\infty=\frac{499}{10000}.
\]
At the root the symmetric packing gate is
$2A_0\log_2r_0<1/10$; below it the old-side asymmetric gate is
$A_d\log_2r_d<1/10$. The exact integer inequality
$6^{200}<2^{517}$ gives $\log_2 6<517/200$. Using this bound for six,
the margins at the five alphabet regimes are at least
\[
 \frac1{12500},\quad\frac1{12500},\quad\frac1{10000},\quad
 \frac{359}{4000000},\quad\frac1{5000},
\]
respectively. The other packing exponents are bounded by
$2/10+A_\infty=2499/10000<1$. Newly created degrees have exponent
at most $A_\infty+o(1)<9/10$. These are strict fixed margins.

Here is a noncircular depth justification. Verify the gates first for
hypothetical active depths through
$\lceil\log_{5/3}\lambda\rceil+O(1)$. Throughout that range,
\[
 u_d\ge\lambda^{\epsilon+o(1)},\qquad
 \epsilon=1-\frac{\ln(8/5)}{\ln(5/3)}>0,
\]
since $8/5<5/3$. Thus floors are negligible, and the total buffers cost
only $O((\log\lambda)^2)=o(\lambda)$. The strict gates imply nearly
lossless padding, and eventually every child has at most $3M/5$ wires.
Consequently no active path extends beyond the proposed range. The
sharper shrink ratios give, uniformly on active paths,
\[
 M_d\le(1+o(1))s f_d,
 \quad(f_0,f_1,f_2,f_3)=\left(1,\frac18,\frac1{64},\frac1{256}\right),
 \quad f_{4+k}=\frac1{768}\,2^{-k}.
\]
The inherited precision bound is $O(s\lambda^2 2^{a_{d-1}})$.
The buffer gives $D\ge2^{a_{d-1}+b+1}$, so its ratio to $MD/u$ is
$O(\lambda^{-2})$ while $M>H$. This establishes the sharp precision
and entropy premises uniformly, including at the final active depths.

\begin{theorem}[Polynomial-space adaptive-quartet certificate]
\label{thm:quartet-final}
Under the normalized source contraction interface, adaptive quartet
conditioning with~\eqref{eq:quartet-schedule} takes
\[
 (s+\mathcal L)^{O(1)}\,
 2^{(C_{\rm quartet}+o(1))s/\sqrt{\lceil\log_2s\rceil}}
\]
time and polynomial space, where
\begin{align*}
 C_{\rm quartet}=\kappa\bigg[&
 \frac{9\sqrt{442/3}}{128\sqrt{1249/100000}}
 +\frac{\sqrt{110/3}}{512\sqrt{26/3125}}
 +\frac{\sqrt{22}}{1536\sqrt{107/20000}}\\
 &+\frac{\sqrt{26/3}}
 {3072\sqrt{27/6400}\,(1-\sqrt{2/5})}\bigg],
\end{align*}
with the certified interval
\[
 50.30468921409<C_{\rm quartet}<50.30468921410.
\]
In particular, $50.305$ is a valid eventual coefficient.
\end{theorem}
\begin{proof}
Sum the logarithms of the branching factors in
\eqref{eq:quartet-local} along a path. The tail has ratio
$(1/2)/\sqrt{5/8}=\sqrt{2/5}$; substituting the ledger gives the
displayed expression. The $O(\log\lambda)$ levels of polynomial
factors contribute $O(\lambda\log\lambda)$ to the logarithm of time.
Leaves cost $2^{O(s/\lambda)}$, initial entry routines cost
$2^{O(s^{9/10})}$, and freshly created packet routines cost
$\exp(s^{499/10000+o(1)})$. Each is lower order than
$s/\sqrt\lambda$, including after paying ancestor child-call factors.
Canonical creation-time descriptions prevent repeated expansion of old
entry routines. Depth-first traversal proves the space claim. The
coefficient interval follows by exact rational arithmetic in the
accompanying checker.
\end{proof}

For a particularly simple structural comparison, keep the incoming
adaptive packet's fixed $r=16$ schedule unchanged. Replacing its balanced
binary trees by adaptive quartets gives
\[
 C_{16}=\frac{5\pi}{2\ln2}\sqrt{\frac{221}{15}}
 \left(1+\frac1{3\sqrt3}\right)
 =51.862701863893\ldots,
\]
exactly $\sqrt{13/15}$ times its previous $55.70947831891\ldots$
coefficient. The improved strict schedule then lowers this to
$50.30468921410$.

\subsection{Comparison and verification boundary}

The two compatible incoming methods already compose to a stronger
polynomial-space comparator than either headline alone: combine the
heterogeneous $(13,14,8,6,4,\ldots)$ schedule in the frontier packet with
the adaptive packet's fixed binary separators. This gives
$55.40466590661\ldots$. A numerical search suggests a fixed-tree
$(16,16,8,6,4,\ldots)$ boundary of $54.00208\ldots$, but it is not
needed for any proof or global optimality claim. At the same strict
rational schedule~\eqref{eq:quartet-schedule}, the old fixed trees have
coefficient $54.02470271284\ldots$, whereas the new quartet trees have
$50.30468921410\ldots$. This matched comparison isolates the new
structural gain from schedule retuning.

The checker performs $15{,}625$ exhaustive four-group inequalities,
$2{,}000$ exact integer-weight profile checks on $r=4,6,8,16$, and
$80$ literal-versus-streamed contractions over
$(\mathbb Z/64\mathbb Z)[V]/(V^2)$. It also verifies every individual
leaf-call count ($14{,}056$ calls in total), the unequal-block norm
ledgers, strict capacity inequalities, and rational coefficient bounds.
Three matched fixtures additionally reduce leaf-factor calls from $64$ to $16$,
$128$ to $32$, and $256$ to $64$ for four, eight, and sixteen groups,
respectively, with exactly equal outputs; these are operation counts,
not timings. All tests remain active under optimized Python. These finite checks
support the implementation and numerical certificates; the all-size
statement rests on the proofs above and the explicitly imported source
interfaces. There is no new proof-assistant build or complete SAT engine.
The result improves a normalized-wire exponent coefficient and preserves
the polynomial-space guarantee; the SAT logarithmic order is unchanged.
